% talk a little about initial page table conditions:
%     paging not on, but virtual mostly mapped direct to physical,
%     which is what things look like when we turn paging on as well
%     since paging is turned on after we create first process.
%   mention why still have SEG_UCODE/SEG_UDATA?
%   do we ever really say what the low two bits of %cs do?
%     in particular their interaction with PTE_U
%   sidebar about why it is extern char[]
\chapter{页表}
\label{CH:MEM}

页表（page table）是操作系统为每个进程提供私有地址空间和内存的最流行机制。页表决定内存地址的含义，以及哪些物理内存区域可被访问。它们允许 xv6 隔离不同进程的地址空间，并将其多路复用到单个物理内存上。页表提供了一层间接性，使操作系统能够执行许多有用的技巧。Xv6 实现了一些：在多个地址空间中映射相同的内存（trampoline 页），用未映射的页保护内核和用户栈，以及延迟分配用户堆内存。本章的其余部分将解释 RISC-V 硬件提供的页表以及 xv6 如何使用它们。

%%
\section{分页硬件}
%%
提醒一下，RISC-V 指令（用户态和内核态）操作虚拟地址（virtual address）。机器的 RAM（即物理内存，physical memory）使用物理地址索引。RISC-V 页表硬件通过将每个虚拟地址映射到物理地址来连接这两种地址。

\begin{figure}[t]
\centering
\includegraphics[scale=0.5]{fig/riscv_address.pdf}
\caption{平坦页表将虚拟地址映射到物理地址的抽象视图。}
\label{fig:riscv_address}
\end{figure}

Xv6 使用 RISC-V 的 Sv39 模式，这意味着 64 位虚拟地址只有低 39 位被使用，高 25 位不使用。在此 Sv39 配置中，RISC-V 页表逻辑上是 $2^{27}$ (134,217,728) 个 \indextext{page table entries (PTEs)}（页表项）的数组。每个 PTE 包含一个 44 位的物理页号（PPN）和一些标志。分页硬件使用 39 位中的高 27 位索引到页表中找到 PTE，并生成一个 56 位的物理地址，其高 44 位来自 PTE 中的 PPN，低 12 位从原始虚拟地址复制。
Figure~\ref{fig:riscv_address}
显示了此过程，页表的逻辑视图是 PTE 的简单数组（RISC-V 页表实际上是一棵树；参见 Figure~\ref{fig:riscv_pagetable} 了解完整故事）。页表使操作系统能够以 4096 ($2^{12}$) 字节对齐块的粒度控制虚拟到物理地址转换。这样的块称为 \indextext{page}（页）。

\begin{figure}[t]
\centering
\includegraphics[scale=0.5]{fig/riscv_pagetable.pdf}
\caption{RISC-V 地址转换细节。}
\label{fig:riscv_pagetable}
\end{figure}

RISC-V 的设计为虚拟地址和物理地址的扩展留下了空间。如果需要更多虚拟地址空间，RISC-V 支持 Sv48 模式，具有 48 位虚拟地址~\cite{riscv:priv}。物理地址也有扩展空间：PTE 格式中有空间让物理页号再增长 10 位。RISC-V 的设计者根据技术预测选择了地址大小。$2^{48}$ 字节即 262,144 GB，远大于当今任何应用程序可能使用的用户虚拟地址空间。$2^{56}$ 字节即 65,536 TB 的物理地址空间，远超目前任何计算机可配备的 RAM。

如
Figure~\ref{fig:riscv_pagetable}
所示，RISC-V CPU 页表作为三级树存储在物理内存中。树的根是一个 4096 字节的页表页，包含 512 个 PTE，其中包含树下一级页表页的物理地址。这些页中的每一个都包含树最后一级的 512 个 PTE。分页硬件使用 27 位中的高 9 位在根页表页中选择 PTE，中间 9 位在树下一级的页表页中选择 PTE，低 9 位选择最终的 PTE。（在 Sv48 RISC-V 中，页表有四级，虚拟地址的 39 到 47 位索引到顶级。）

如果转换地址所需的三个 PTE 中任何一个不存在，分页硬件会引发 \indextext{page-fault exception}（页错误异常），由内核处理页错误（参见 Chapters~\ref{CH:TRAP} 和~\ref{CH:PGFAULTS}）。

与 Figure~\ref{fig:riscv_address} 的单级设计相比，Figure~\ref{fig:riscv_pagetable} 的三级结构以更节省内存的方式记录 PTE。在虚拟地址大范围无映射的常见情况下，三级结构可以省略整个页目录。例如，如果应用程序只使用从地址零开始的少数几页，那么顶级页目录的条目 1 到 511 无效，内核无需为这 511 个中间页目录分配页面。此外，内核也无需为这 511 个中间页目录的底层页目录分配页面。因此，在此例中，三级设计节省了 511 个中间页目录页面和 $511\times512$ 个底层页目录页面。

尽管 CPU 在执行加载或存储指令时会以硬件方式遍历三级结构，但三级的一个潜在缺点是 CPU 必须从内存加载三个 PTE 才能将加载/存储指令中的虚拟地址转换为物理地址。为了避免从物理内存加载 PTE 的开销，RISC-V CPU 将页表项缓存在
\indextext{Translation Look-aside Buffer (TLB)}（转换后备缓冲器）中。

每个 PTE 包含标志位，告知分页硬件如何使用关联的虚拟地址。
\indexcode{PTE_V}
指示 PTE 是否存在：若未设置，对页面的引用会导致页错误（即不允许）。
\indexcode{PTE_R}
控制是否允许指令读取页面。
\indexcode{PTE_W}
控制是否允许指令写入页面。
\indexcode{PTE_X}
控制 CPU 是否可将页面内容解释为指令并执行。
\indexcode{PTE_U}
控制用户模式中的指令是否允许访问页面；若
\indexcode{PTE_U}
未设置，PTE 只能在监管者模式中使用。
Figure~\ref{fig:riscv_pagetable}
显示了标志位在 PTE 中的位置。标志和所有其他与页面硬件相关的结构定义在
\fileref{kernel/riscv.h}
中。

为了让 CPU 使用页表，内核必须将根页表页的物理地址写入
\texttt{satp}\index{satp@\lstinline{satp}}
寄存器。CPU 将使用其 \texttt{satp} 指向的页表翻译后续指令生成的所有地址。每个 CPU 都有自己的 \texttt{satp}，因此不同的 CPU 可以运行不同的进程，每个进程都有自己的页表描述的私有地址空间。

从内核的角度看，页表是存储在内存中的数据，内核使用类似于树形数据结构的代码创建和修改页表。

关于本书中使用的一些术语的说明。\textit{物理内存}（Physical memory）指 RAM 中的存储单元。物理内存的字节有一个地址，称为 \textit{物理地址}（physical address）。解引用地址的指令（如加载、存储、跳转和函数调用）只使用虚拟地址（virtual address），分页硬件将其转换为物理地址，然后发送到 RAM 硬件读取或写入存储。\textit{地址空间}（address space）是给定页表中有效的虚拟地址集合；每个 xv6 进程有独立的用户地址空间，xv6 内核也有自己的地址空间。\textit{用户内存}（User memory）指进程的用户地址空间加上页表允许进程访问的物理内存。\textit{虚拟内存}（Virtual memory）指与页表管理及使用它们实现隔离等相关的概念和技术。

\begin{figure}[h]
\centering
 \includegraphics[scale=0.65]{fig/xv6_layout.pdf}
\caption{左侧，xv6 的内核虚拟地址空间。
{\sf \small{RWX}}
指 PTE 的读、写和执行权限。
右侧，xv6 期望看到的 RISC-V 物理地址空间。}
\label{fig:xv6_layout}
\end{figure}

%%
\section{内核地址空间}
%%
启动时，xv6 创建一个描述内核地址空间的页表。内核配置其地址空间布局，使自己能够访问物理内存和各种硬件资源的可预测虚拟地址。
Figure~\ref{fig:xv6_layout}
显示了此布局如何将内核虚拟地址映射到物理地址。文件
\fileref{kernel/memlayout.h}
声明了 xv6 内核内存布局的常量。

QEMU 模拟一台计算机，包括从物理地址 \texttt{0x80000000} 开始并至少持续到 \texttt{0x88000000} 的 RAM（物理内存），xv6 称之为 \texttt{PHYSTOP}。QEMU 模拟还包括 I/O 设备，如磁盘接口。QEMU 将设备接口作为位于物理地址空间 \texttt{0x80000000} 以下的 \indextext{memory-mapped}（内存映射）控制寄存器暴露给软件。内核可以通过读写这些特殊物理地址与设备交互；此类读写与设备硬件通信而非 RAM。Chapter~\ref{CH:TRAP} 解释 xv6 如何与设备交互。

内核将所有物理 RAM 和设备寄存器映射到等于物理地址的虚拟地址。这称为``直接映射''，允许内核通过简单地加载或存储虚拟地址 $x$ 来读取或写入物理地址 $x$。内核代码本身位于虚拟地址空间和物理内存中的 \lstinline{KERNBASE=0x80000000}。当 \lstinline{kfork}
\lineref{kernel/proc.c:/^kfork/}
为子进程分配用户内存时，分配器返回该内存的物理地址；\lstinline{kfork} 在将父进程的用户内存复制到子进程时直接使用该地址作为虚拟地址。

有几个内核虚拟地址不是直接映射的：

\begin{itemize}
  
\item trampoline 页。它映射在虚拟地址空间的顶部；用户页表具有相同的映射。Chapter~\ref{CH:TRAP} 讨论了 trampoline 页的作用，但我们在这里看到了页表的一个有趣用例：一个物理页（保存 trampoline 代码）在内核的虚拟地址空间中映射了两次：一次在虚拟地址空间的顶部，一次使用直接映射。

\item 内核栈页。每个进程都有自己的内核栈，映射在较高的内核虚拟地址处，因此在其下方 xv6 可以保留一个未映射的\indextext{保护页}。保护页的 PTE 是无效的（即\lstinline{PTE_V}未设置），因此如果内核溢出内核栈，很可能会引发页故障，导致内核崩溃。如果没有保护页，溢出的栈将覆盖其他内核内存，导致错误操作。崩溃重启是可取的。

\end{itemize}

虽然内核通过高内存映射使用其栈，但每个栈也可通过直接映射地址供内核访问。另一种设计可能只有直接映射，并在直接映射地址使用栈。然而，在这种安排中，提供保护页将涉及取消映射否则会引用物理内存的虚拟地址，这将很难使用。

内核将 trampoline 和内核文本页映射为具有 \lstinline{PTE_R} 和 \lstinline{PTE_X} 权限，但没有 \lstinline{PTE_W}。内核将其他页映射为具有 \lstinline{PTE_R} 和 \lstinline{PTE_W} 权限，但没有 \lstinline{PTE_X}。保护页映射是无效的。这些受限权限旨在帮助捕获以意外方式访问页的内核错误，例如若内核代码意外尝试写入内核指令。

内核创建一个单一的内核页表，所有 CPU 在内核中执行时使用它。xv6 在最初创建后不会修改内核页表。

%%
\section{代码：创建地址空间}
%%

请在继续之前阅读 {\tt kernel/vm.c} 到 {\tt mappages()} 的结尾。

大部分用于操作地址空间和页表的 xv6 代码位于 {\tt vm.c}
\fileref{kernel/vm.c}。
中心数据结构是 {\tt pagetable\_t}，它实际上是指向 RISC-V 根页表页的指针；{\tt pagetable\_t} 可以是内核页表或每个进程的页表之一。中心函数是 {\tt walk}，它查找虚拟地址的 PTE，以及 {\tt mappages}，它安装新映射的 PTE。以 {\tt kvm} 开头的函数操作内核页表；以 {\tt uvm} 开头的函数操作用户页表；其他函数用于两者。{\tt copyout} 和 {\tt copyin} 将数据复制到作为系统调用参数提供的用户虚拟地址或从该地址复制；它们在 {\tt vm.c} 中，因为它们需要显式转换这些地址以找到相应的物理内存。

在启动序列的早期，
\indexcode{main}
调用
\indexcode{kvminit}
\lineref{kernel/vm.c:/^kvminit/}
使用
\indexcode{kvmmake}
\lineref{kernel/vm.c:/^kvmmake/}
创建内核的页表。此调用发生在 xv6 在 RISC-V 上启用分页之前，因此地址直接引用物理内存。
\lstinline{kvmmake}
首先分配一页物理内存来保存根页表页。然后它调用
\indexcode{kvmmap}
来安装内核需要的转换。转换包括内核的指令和数据、高达
\indexcode{PHYSTOP}
的物理内存，以及实际上是设备的内存范围。
\indexcode{proc\_mapstacks}
\lineref{kernel/proc.c:/^proc\_mapstacks/}
为每个进程分配一个内核栈。它调用 \lstinline{kvmmap} 将每个栈映射到由
\lstinline{KSTACK}
生成的虚拟地址，这为无效的栈保护页留出了空间。

\indexcode{kvmmap} 
\lineref{kernel/vm.c:/^kvmmap/}
calls
\indexcode{mappages}
\lineref{kernel/vm.c:/^mappages/}，
它在页表中为一组虚拟地址安装映射到相应的物理地址范围。它以页间隔对范围内的每个虚拟地址执行此操作。对于要映射的每个虚拟地址，
\lstinline{mappages}
调用
\indexcode{walk}
来查找该地址的 PTE 地址。然后它初始化 PTE 以保存相关的物理页号、所需的权限
（\lstinline{PTE_W}、
\lstinline{PTE_X}
和/或
\lstinline{PTE_R}），
以及
\lstinline{PTE_V}
以将 PTE 标记为有效
\lineref{kernel/vm.c:/perm...PTE_V/}。

\indexcode{walk}
\lineref{kernel/vm.c:/^walk/}
模仿 RISC-V 分页硬件查找虚拟地址的 PTE（参见
Figure~\ref{fig:riscv_pagetable}）。
\lstinline{walk}
逐级下降页表树，使用每级的 9 位虚拟地址索引到相关的页目录页。在每一级，它找到下一级页目录页的 PTE 或最终页的 PTE
\lineref{kernel/vm.c:/pte.=..pagetable/}。
如果第一级或第二级页目录页中的 PTE 无效，则所需的目录页尚未分配；如果
\lstinline{alloc}
参数已设置，
\lstinline{walk}
分配一个新的页表页并将其物理地址放入 PTE。它返回树中最底层 PTE 的地址
\lineref{kernel/vm.c:/return..pagetable/}。

上述代码依赖于物理内存直接映射到内核虚拟地址空间。例如，当 \lstinline{walk} 逐级下降页表时，它从 PTE 中提取下一级页表的（物理）地址 \lineref{kernel/vm.c:/pagetable.=..pa.*E2P/}，然后使用该地址作为虚拟地址来获取下一级 PTE
\lineref{kernel/vm.c:/t..pte.=..paget/}。

在每个 CPU 上，
\indexcode{main}
调用
\indexcode{kvminithart}
\lineref{kernel/vm.c:/^kvminithart/}
来安装内核页表，将根页表页的物理地址放入 CPU 的
\texttt{satp}
寄存器。此后，CPU 使用内核页表翻译地址。内核继续正确执行，因为内核页表是直接映射的，因此地址在此更改前后引用 RAM 中的相同位置。

每个 RISC-V CPU 在
\indextext{Translation Look-aside Buffer (TLB)}（转换后备缓冲器）中缓存页表项，当 xv6 更改页表时，它必须告诉 CPU 使相应的缓存 TLB 条目无效。如果没有这样做，那么在之后的某个时刻，TLB 可能会使用旧的缓存映射，指向在此期间已分配给另一个进程的物理页，结果进程可能会在其他进程的内存上涂鸦。RISC-V 有一条指令 \indexcode{sfence.vma} 刷新当前 CPU 的 TLB。Xv6 在重新加载
\texttt{satp}
寄存器后在 {\tt kvminithart} 中执行 {\tt sfence.vma}，以及在
\lstinline{uservec}
和 \lstinline{userret} 的 trampoline 代码中执行。

在更改
\texttt{satp}
之前，还需要发出 \texttt{sfence.vma}，以等待所有未完成的加载和存储完成。此等待确保对页表的先前更新已完成，并确保先前的加载和存储使用旧页表，而不是新页表。

\section{物理内存分配}

内核必须在运行时分配和释放物理内存，用于页表、用户内存、内核栈和管道缓冲区。

Xv6 使用内核末尾和 \indexcode{PHYSTOP} 之间的物理内存进行运行时分配。它每次分配和释放整个 4096 字节的页。它通过在页面本身中内嵌链表来跟踪哪些页面是空闲的。分配即把一页从链表中删除；释放即将页面添加到链表中。

请阅读 {\tt kernel/kalloc.c}。

%%
\section{代码：物理内存分配器}
%%

分配器位于 {\tt kalloc.c} \fileref{kernel/kalloc.c}。其数据结构是一个 \textit{空闲列表}（free list），包含可用于分配的物理内存页。每个空闲页存储其``下一个''指针的位置在 \indexcode{struct run} \lineref{kernel/kalloc.c:/^struct.run/}。分配器将每个空闲页的 \lstinline{run} 结构存储在空闲页本身中，因为页面空闲时无其他内容存储。空闲列表受自旋锁保护 \linerefs{kernel/kalloc.c:/^struct.{/,/}/}。列表和锁包装在一个结构体中，以明确锁保护结构体中的字段。现在，忽略锁和对 \lstinline{acquire} 与 \lstinline{release} 的调用；Chapter~\ref{CH:LOCK} 将详细讨论锁定。

函数
\indexcode{main}
调用
\indexcode{kinit}
来初始化分配器
\lineref{kernel/kalloc.c:/^kinit/}。
\lstinline{kinit} 将空闲列表初始化为包含内核末尾和 {\tt PHYSTOP} 之间的每个物理 RAM 页。Xv6 本应该通过解析硬件提供的配置信息来确定有多少物理内存可用。但实际上，xv6 假设机器有 128 MB 的 RAM。
\lstinline{kinit}
调用
\indexcode{freerange}
通过每页调用
\indexcode{kfree}
来将内存添加到空闲列表。PTE 只能引用对齐到 4096 字节边界（是 4096 的倍数）的物理地址，因此
\lstinline{freerange}
使用
\indexcode{PGROUNDUP}
来确保它仅释放对齐的物理地址。分配器从没有内存开始；这些对
\lstinline{kfree}
的调用给它一些内存来管理。

分配器有时将地址视为整数以便对它们执行算术运算（例如，在 \lstinline{freerange} 中遍历所有页），有时使用地址作为指针来读写内存（例如，操作存储在每页中的 \lstinline{run} 结构）；地址的这种双重用途是分配器代码充满 C 类型转换的主要原因。
\index{type cast}

函数
\lstinline{kfree}
\lineref{kernel/kalloc.c:/^kfree/}
首先将正在释放的内存中的每个字节设置为值 1。这将导致在释放内存后使用内存的代码（使用``悬空引用''）读取垃圾而不是旧的 valid 内容；希望这将导致此类代码更快地崩溃。然后
\lstinline{kfree}
将页面前置到空闲列表：它将
\lstinline{pa}
转换为指向
\lstinline{struct}
\lstinline{run}
的指针，将空闲列表的旧开始记录在
\lstinline{r->next}，
并将空闲列表设置为
\lstinline{r}。
\indexcode{kalloc}
从空闲列表中移除并返回第一个元素。

\section{进程地址空间}

每个进程都有自己的页表，当 xv6 在进程之间切换时，它也会更改页表。
Figure~\ref{fig:processlayout} 比 Figure~\ref{fig:as} 更详细地显示了进程的地址空间。进程的用户地址空间从零开始，理论上结束于
\texttt{MAXVA}
({\tt 0x4000000000})
\lineref{kernel/riscv.h:/MAXVA/}，尽管实际上只有其中一小部分映射到物理内存。

进程的地址空间由以下页面组成：包含程序文本的页面（xv6 使用权限 \lstinline{PTE_R}、\lstinline{PTE_X} 和 \lstinline{PTE_U} 映射）、包含程序预初始化数据的页面、栈页面以及堆页面。Xv6 使用权限 \lstinline{PTE_R}、\lstinline{PTE_W} 和 \lstinline{PTE_U} 映射数据、栈和堆。

在用户地址空间内使用权限是强化用户进程的常用技术。如果文本使用 \lstinline{PTE_W} 映射，那么进程可能会意外修改自己的程序；例如，编程错误可能导致程序写入空指针，修改地址 0 处的指令，然后继续运行，可能造成更大的破坏。为了立即检测此类错误，xv6 在不使用 \lstinline{PTE_W} 的情况下映射文本；如果程序意外尝试存储到地址 0，硬件将拒绝执行存储并引发页错误（参见 Chapter~\ref{CH:TRAP}）。然后内核杀死进程并打印出信息性消息，以便开发人员可以追踪问题。

同样，通过在不使用 \lstinline{PTE_X} 的情况下映射数据，用户程序不能意外地跳转到程序数据中的地址并从该地址开始执行。

在现实世界中，通过仔细设置权限来强化进程也有助于防御安全攻击。攻击者可能向程序（例如 Web 服务器）提供精心构造的输入，触发程序中的错误，希望将该错误转化为漏洞利用~\cite{aleph:smashing}。仔细设置权限和其他技术（如用户地址空间布局随机化）使此类攻击更难。

栈是单页，显示为 {\tt exec} 系统调用创建的初始内容。包含命令行参数的字符串以及指向它们的指针数组位于栈的最顶部。就在那下面是允许程序在
\lstinline{main}
开始执行，就好像函数 \lstinline{main(argc}、\lstinline{argv)} 刚刚被调用一样。

为了检测用户栈溢出分配的栈内存，xv6 通过清除 \lstinline{PTE_U} 标志在栈正下方放置一个不可访问的保护页。如果用户栈溢出并且进程尝试使用栈下方的地址，硬件将生成页错误异常，因为保护页对用户模式下运行的程序不可访问。现实世界的操作系统可能会在栈溢出时自动为用户栈分配更多内存。

我们在这里看到了一些页表使用的很好的例子。首先，不同进程的页表将用户地址转换为不同的物理内存页，因此每个进程都有私有的用户内存。其次，每个进程将其内存视为具有从零开始的连续虚拟地址，而进程的物理内存可以是非连续的。第三，内核在用户地址空间顶部映射一个具有 trampoline 代码的页（没有 \lstinline{PTE_U}），因此单个物理内存页出现在所有地址空间中，但只能由内核使用。

\begin{figure}[t]
\centering
\includegraphics[scale=0.5]{fig/processlayout.png}
\caption{进程的用户地址空间及其初始栈。}
\label{fig:processlayout}
\end{figure}


%%
\section{代码：exec}
%%

请阅读 {\tt kernel/exec.c} 和 {\tt kernel/vm.c}，从 {\tt uvmcreate()} 开始。

\lstinline{exec} 是一个系统调用，用从文件读取的数据（称为二进制或可执行文件）替换进程的用户地址空间。二进制通常是编译器和链接器的输出，保存机器指令和程序数据。
\lstinline{kexec} \lineref{kernel/exec.c:/^kexec/}，
{\tt exec} 的内核内部实现，使用 \indexcode{namei}
\lineref{kernel/exec.c:/namei/}
打开命名的二进制文件 \lstinline{path}，这在 Chapter~\ref{CH:FS} 中解释。然后，它读取 ELF 头。Xv6 二进制文件以广泛使用的 \indextext{ELF format}（ELF 格式）格式化，定义在
\fileref{kernel/elf.h}。ELF 二进制文件由 ELF 头
\indexcode{struct elfhdr} \lineref{kernel/elf.h:/^struct.elfhdr/}、
后跟一系列程序段头
\lstinline{struct proghdr} \lineref{kernel/elf.h:/^struct.proghdr/} 组成。每个
\lstinline{proghdr}
描述必须加载到内存中的应用段；xv6 程序有两个程序段头：一个用于指令，一个用于数据。

第一步是快速检查文件是否可能包含 ELF 二进制文件。ELF 二进制文件以四字节``魔数''
\lstinline{0x7F}、
\lstinline{`E'}、
\lstinline{`L'}、
\lstinline{`F'}
或
\indexcode{ELF_MAGIC}
\lineref{kernel/elf.h:/ELF_MAGIC/} 开头。
如果 ELF 头具有正确的魔数，
\lstinline{kexec}
假设二进制文件格式正确。

\lstinline{kexec}
使用
\indexcode{proc\_pagetable}
\lineref{kernel/exec.c:/proc\_pagetable/}
分配一个没有用户映射的新页表，使用
\indexcode{uvmalloc}
\lineref{kernel/exec.c:/uvmalloc/}
为每个 ELF 段分配内存，并使用
\indexcode{loadseg}
\lineref{kernel/exec.c:/loadseg/}
将每个段加载到内存中。
\lstinline{loadseg}
使用
\indexcode{walkaddr}
来查找分配的内存的物理地址，以便写入
ELF 段的每一页，并使用
\indexcode{readi}
从文件中读取。

\indexcode{/init}
的程序段头，
用
\lstinline{exec}
创建的第一个用户程序，
看起来像这样：
\begin{footnotesize}
\begin{verbatim}
# objdump -p user/_init

user/_init:     file format elf64-little

Program Header:
0x70000003 off    0x0000000000006bb0 vaddr 0x0000000000000000
                                       paddr 0x0000000000000000 align 2**0
         filesz 0x000000000000004a memsz 0x0000000000000000 flags r--
    LOAD off    0x0000000000001000 vaddr 0x0000000000000000
                                       paddr 0x0000000000000000 align 2**12
         filesz 0x0000000000001000 memsz 0x0000000000001000 flags r-x
    LOAD off    0x0000000000002000 vaddr 0x0000000000001000
                                       paddr 0x0000000000001000 align 2**12
         filesz 0x0000000000000010 memsz 0x0000000000000030 flags rw-
   STACK off    0x0000000000000000 vaddr 0x0000000000000000
                                       paddr 0x0000000000000000 align 2**4
         filesz 0x0000000000000000 memsz 0x0000000000000000 flags rw-
\end{verbatim}
\end{footnotesize}

我们看到文本段应该加载到内存中的虚拟地址 0x0（无写权限），从文件中偏移量 0x1000 处的内容加载。我们还看到数据应该加载到地址 0x1000，该地址位于页边界，并且没有可执行权限。

程序段头的
\lstinline{filesz}
可能小于
\lstinline{memsz}，
表示它们之间的间隙应该用零填充（对于 C 全局变量）而不是从文件中读取。
对于
\lstinline{/init}，数据
\lstinline{filesz}
为 0x10 字节，
\lstinline{memsz}
为 0x30 字节，
因此
\indexcode{uvmalloc}
分配足够的物理内存以容纳 0x30 字节，但仅从文件读取 0x10 字节
\lstinline{/init}。

现在
\indexcode{kexec}
分配并初始化用户栈。它只分配一个栈页。
\lstinline{kexec}
将参数字符串一次一个地复制到栈顶，在
\indexcode{ustack}
中记录指向它们的指针。它在
\indexcode{argv}
列表的末尾放置一个空指针，该列表将传递给
\lstinline{main}。
\indexcode{argc}
和 \lstinline{argv} 的值通过系统调用返回路径传递给 \lstinline{main}：\lstinline{argc} 通过系统调用返回值传递，放在 {\tt a0} 中，\lstinline{argv} 通过进程 trapframe 的 {\tt a1} 条目传递。

\lstinline{kexec}
在栈页正下方放置一个不可访问的页，因此尝试使用多个页的程序将出错。此不可访问的页还允许
\lstinline{kexec}
处理太大的参数；在这种情况下，
\lstinline{kexec}
使用的
\indexcode{copyout}
\lineref{kernel/vm.c:/^copyout/}
函数将注意到目标页不可访问，并将返回 -1。

在准备新内存映像期间，如果
\lstinline{kexec}
检测到错误（如无效程序段），它会跳转到标签
\lstinline{bad}，
释放新映像，并返回 -1。
\lstinline{kexec}
必须等到确定系统调用成功后才能释放旧映像：如果旧映像已消失，系统调用无法向它返回 -1。
\lstinline{kexec}
中的唯一错误情况发生在映像创建期间。一旦映像完成，
\lstinline{kexec}
就可以提交到新页表
\lineref{kernel/exec.c:/pagetable.=.pagetable/}
并释放旧页表
\lineref{kernel/exec.c:/proc_freepagetable/}。

\lstinline{exec} 系统调用将字节从 ELF 文件加载到 ELF 文件指定的内存地址。用户或进程可以在 ELF 文件中放置任意地址。因此 \lstinline{exec} 是有风险的，因为 ELF 文件中的地址可能意外或故意引用内核。对粗心的内核而言，后果可能从崩溃到恶意颠覆内核的隔离机制（即安全漏洞）。Xv6 执行许多检查以避免这些风险。例如，\lstinline{if(ph.vaddr + ph.memsz < ph.vaddr)} 检查总和是否溢出 64 位整数。危险在于用户可以构造一个 ELF 二进制文件，其 \lstinline{ph.vaddr} 指向用户选择的地址，且 \lstinline{ph.memsz} 足够大，使总和溢出到 0x1000，这看起来像一个有效值。在 xv6 的旧版本中，用户地址空间也包含内核（但在用户模式下不可读/写），用户可以选择对应于内核内存的地址，从而将数据从 ELF 二进制文件复制到内核中。在 RISC-V 版本的 xv6 中，这不可能发生，因为内核有自己的独立页表；
\lstinline{loadseg}
加载到进程的页表中，而不是内核的页表中。

内核开发人员很容易省略关键检查，现实世界的内核长期缺少可以被用户程序利用来获取内核权限的检查。xv6 可能未完全验证提供给内核的用户级数据，恶意用户程序可能利用这一点来绕过 xv6 的隔离。
%%
\section{现实世界}
%%

与大多数操作系统一样，xv6 使用分页硬件进行内存保护和映射。大多数操作系统通过结合分页和页错误异常更复杂地使用分页，我们将在 Chapter~\ref{CH:TRAP} 中讨论。

Xv6 因内核使用虚拟地址和物理地址之间的直接映射，以及假设物理 RAM 位于地址 0x80000000（内核期望加载的位置）而简化。这适用于 QEMU，但在真实硬件上效果不佳；真实硬件将 RAM 和设备放置在不可预测的物理地址，因此（例如）0x80000000 处可能没有 RAM，而 xv6 期望能够在该处存储内核。更严肃的内核设计利用页表将任意的硬件物理内存布局转换为可预测的内核虚拟地址布局。

RISC-V 支持物理地址级别的保护，但 xv6 不使用该功能。

在内存较多的机器上，使用 RISC-V 对``超级页''的支持可能是有意义的。小页面在物理内存较小时有意义，以允许细粒度分配和换出到磁盘。例如，如果程序仅使用 8 KB 内存，为其分配整个 4 MB 的超级物理内存页是浪费的。较大的页面在 RAM 较多的机器上有意义，并可能减少页表操作的开销。

为了避免在更改页表时必须刷新整个 TLB，RISC-V CPU 可能支持地址空间标识符（ASIDs）~\cite{riscv:priv}。然后内核可以只刷新特定地址空间的 TLB 条目。Xv6 不使用此功能。

xv6 内核缺乏类似 {\tt malloc} 的分配器来为小对象提供内存，这阻止了内核使用需要动态分配的复杂数据结构。更复杂的内核可能会分配许多不同大小的小块，而不是（如 xv6）仅 4096 字节块；真正的内核分配器需要处理小分配和大分配。

内存分配是一个永恒的热门话题，基本问题是有效使用有限内存和为未知的未来请求做准备~\cite{knuth}。今天人们更关心速度而不是空间效率。
%%
\section{练习}
%%

\begin{enumerate}

\item 解析 RISC-V 的设备树以查找计算机拥有的物理内存量。

\item 函数 {\tt copyin} 和 {\tt copyinstr} 在软件中遍历用户页表。设置内核页表，使内核映射了用户程序，并且 {\tt copyin} 和 {\tt copyinstr} 可以使用 {\tt memcpy} 将系统调用参数复制到内核空间，依靠硬件进行页表遍历。

\item 修改 xv6 以在内核中使用超级页。

\item Unix 的
\lstinline{exec}
实现对 shell 脚本有特殊处理。如果要执行的文件以文本
\lstinline{#!}
开头，则第一行被视为运行来解释文件的程序。例如，如果
\lstinline{exec}
被调用以运行
\lstinline{myprog}
\lstinline{arg1}
并且
\lstinline{myprog}
的第一行是
\lstinline{#!/interp}，
则
\lstinline{exec}
运行
\lstinline{/interp}
，命令行为
\lstinline{/interp}
\lstinline{myprog}
\lstinline{arg1}。
在 xv6 中实现对此约定的支持。

\item 为内核实现地址空间布局随机化。

\end{enumerate}
